% -------------------------------------------------------------------------------------
\section{QCD 拉氏量及其规范不变性}\label{sec:prelim-qcd}

\cref{prop:qcd-covariance}告诉我们哪些组合在规范变换下协变：$D_\mu\psi$ 与 $\psi$ 一样变换，$G_{\mu\nu}$ 按伴随表示变换。由它们构造标量就得到规范不变的拉氏量。全书使用的拉氏量如下；其中后两项（规范固定项与鬼场项）在经典层次上并不需要，它们的来源与作用在\cref{sec:qcd-quantization}中说明。

\begin{definition}[QCD 拉氏量与规范变换]\label{def:prelim-qcd}
含 $n_f$ 味夸克（质量 $m_q\ge0$）的协变规范 QCD 拉氏量为 % VF-CONV-2
\begin{equation}\label{eq:qcd-lagrangian}
\mathcal L=-\frac14G^a_{\mu\nu}G^{a\mu\nu}+\sum_q\bar\psi_q\big(\ii\sla D-m_q\big)\psi_q-\frac1{2\xi_g}\big(\partial^\mu A^a_\mu\big)^2+\bar c^a\big(-\partial^\mu D^{ac}_\mu\big)c^c ,
\end{equation}
其中作用于夸克的协变导数为 $D_\mu=\partial_\mu-\ii gA^a_\mu t^a$，作用于伴随表示场的协变导数为 $(D_\mu X)^a=D^{ac}_\mu X^c=\partial_\mu X^a+gf^{abc}A^b_\mu X^c$，$G^a_{\mu\nu}=\partial_\mu A^a_\nu-\partial_\nu A^a_\mu+gf^{abc}A^b_\mu A^c_\nu$；$\psi_q$ 是带基础表示色指标的 Grassmann 奇 Dirac 旋量场，$A^a_\mu$ 是实的胶子场，$c^a$、$\bar c^a$ 是伴随表示中彼此独立的 Grassmann 奇标量场（鬼场与反鬼场，$\bar c$ 不是 $c$ 的复共轭），$\xi_g$ 是规范参数。规范变换为
\begin{equation}\label{eq:qcd-gauge-trafo}
\psi_q\to U\psi_q,\qquad A_\mu\to UA_\mu U^{-1}-\frac{\ii}g(\partial_\mu U)U^{-1},\qquad U(x)\in\mathrm{SU}(N).
\end{equation}
前两项组成经典拉氏量 $\mathcal L_{\mathrm{cl}}$。约定：Collins\cite{Collins2011}、BBL、LaMET 与 Yizhuang Liu 等人的文献取 $D=\partial+\ii gA$，相对于本书 $g\to-g$。
\uses{def:qcd-gauge-group}
\end{definition}

后两项在\cref{def:qcd-fp}中作为规范固定与 Faddeev--Popov 构造的结果重新给出。由 $G_{\mu\nu}=G^a_{\mu\nu}t^a$ 与 $\Tr(t^at^b)=\frac12\delta^{ab}$，
\begin{equation}\label{eq:qcd-GG-trace}
-\frac14G^a_{\mu\nu}G^{a\mu\nu}=-\frac12\Tr\big(G_{\mu\nu}G^{\mu\nu}\big).
\end{equation}

\begin{theorem}[经典 QCD 拉氏量的规范不变性与唯一性 \tagR]\label{thm:qcd-gauge-invariance}
（$R_{\mathrm{def}}$：由算符定义直接推出。）
\begin{enumerate}[label=(\alph*),leftmargin=2.2em,itemsep=1pt,topsep=2pt]
\item $\mathcal L_{\mathrm{cl}}=-\frac14G^a_{\mu\nu}G^{a\mu\nu}+\sum_q\bar\psi_q(\ii\sla D-m_q)\psi_q$ 在\cref{def:prelim-qcd}的一切局域 $\mathrm{SU}(N)$ 规范变换下不变（有限的 $U(x)$，而不仅是无穷小变换）。项 $G^a_{\mu\nu}\tilde G^{a\mu\nu}$（$\tilde G^{a\mu\nu}=\frac12\epsilon^{\mu\nu\rho\sigma}G^a_{\rho\sigma}$）规范不变，且是一个全导数：
\begin{equation}\label{eq:qcd-chern-simons}
G^a_{\mu\nu}\tilde G^{a\mu\nu}=\partial_\mu J^\mu_{\mathrm{CS}},\qquad
J^\mu_{\mathrm{CS}}=2\,\epsilon^{\mu\nu\rho\sigma}\Big(A^a_\nu\partial_\rho A^a_\sigma+\frac g3f^{abc}A^a_\nu A^b_\rho A^c_\sigma\Big).
\end{equation}
\item 只含 $A^a_\mu$、$\psi_q$、$\bar\psi_q$ 及其导数、质量量纲 $\le4$ 的 Lorentz 不变、规范不变的定域多项式，都是 $G^a_{\mu\nu}G^{a\mu\nu}$、$G^a_{\mu\nu}\tilde G^{a\mu\nu}$、$\bar\psi_q\ii\sla D\psi_{q'}$、$\bar\psi_q\gamma_5\ii\sla D\psi_{q'}$、$\bar\psi_q\psi_{q'}$、$\bar\psi_q\ii\gamma_5\psi_{q'}$（系数为味空间矩阵）的线性组合（不计常数算符 $1$）。经味空间的变换可使动能项取正则形式、质量矩阵取对角且元素实而非负；质量矩阵的整体相位与 $\theta$ 项之间由轴矢反常联系，本书不用到这一点。
\end{enumerate}
（本定理是关于经典场的代数陈述。）
\uses{def:prelim-qcd, prop:qcd-covariance, lem:qcd-colour-identities}
\end{theorem}

\begin{proof}
(a) 由\cref{prop:qcd-covariance}(a)，$G_{\mu\nu}\to UG_{\mu\nu}U^{-1}$，故由\cref{eq:qcd-GG-trace}与迹的循环性，$-\frac12\Tr(G_{\mu\nu}G^{\mu\nu})$ 不变。$U$ 作用于色指标，与 Dirac 矩阵对易，$U^\dagger=U^{-1}$，故 $\bar\psi=\psi^\dagger\gamma^0\to\bar\psi U^{-1}$；又 $\sla D\psi\to U\sla D\psi$，于是 $\bar\psi(\ii\sla D-m)\psi\to\bar\psi U^{-1}U(\ii\sla D-m)\psi$ 不变。同理 $G^a_{\mu\nu}\tilde G^{a\mu\nu}=\epsilon^{\mu\nu\rho\sigma}\Tr(G_{\mu\nu}G_{\rho\sigma})$ 不变。

全导数性质：在 $\epsilon^{\mu\nu\rho\sigma}$ 的反对称化之下 $G_{\mu\nu}\to2\partial_\mu A_\nu-2\ii gA_\mu A_\nu$，故
\[
\epsilon^{\mu\nu\rho\sigma}\Tr(G_{\mu\nu}G_{\rho\sigma})=4\epsilon^{\mu\nu\rho\sigma}\Tr(\partial_\mu A_\nu\partial_\rho A_\sigma)-8\ii g\,\epsilon^{\mu\nu\rho\sigma}\Tr(\partial_\mu A_\nu A_\rho A_\sigma),
\]
其中四个 $A$ 的项因 $\epsilon^{\mu\nu\rho\sigma}\Tr(A_\mu A_\nu A_\rho A_\sigma)=0$（迹的循环移动是四个指标的奇置换）而消失，两个三次项由迹的循环性合并。另一方面，由\cref{lem:qcd-colour-identities}(iv)，$\epsilon^{\nu\rho\sigma\cdots}\Tr(A_\nu A_\rho A_\sigma)$ 只保留 $\frac{\ii}4f^{abc}$ 部分，即 $f^{abc}\epsilon^{\mu\nu\rho\sigma}A^a_\nu A^b_\rho A^c_\sigma=-4\ii\epsilon^{\mu\nu\rho\sigma}\Tr(A_\nu A_\rho A_\sigma)$，于是 $J^\mu_{\mathrm{CS}}=4\epsilon^{\mu\nu\rho\sigma}\Tr(A_\nu\partial_\rho A_\sigma)-\frac{8\ii g}3\epsilon^{\mu\nu\rho\sigma}\Tr(A_\nu A_\rho A_\sigma)$。求散度：$\partial_\mu\partial_\rho$ 的对称性去掉一项，三次项的三个导数位置由循环性与 $\epsilon$ 的偶置换给出相同的贡献，得到上式右边。（这一恒等式也对 $\mathrm{SU}(2)$ 的多项式场作了符号验证。）% s2_lagrangian.py

(b)（证明概要）第一步：规范不变的定域多项式可以写成 $\psi$、$\bar\psi$、$G_{\mu\nu}$ 及它们的协变导数 $D_\lambda\cdots$ 的多项式；这是规范不变性的标准推论（规范变换可以在任一点把 $A_\mu$ 及其对称化导数变到零，剩下的不变量只能通过 $G_{\mu\nu}$ 及其协变导数依赖于 $A$）。第二步：按量纲枚举。$[D_\mu]=1$，$[G_{\mu\nu}]=2$，$[\psi]=\frac32$。纯规范部分：$\Tr G_{\mu\nu}=0$，量纲 $\le3$ 的 Lorentz 标量不存在，量纲 $4$ 的只有 $\Tr(G_{\mu\nu}G^{\mu\nu})$ 与 $\epsilon^{\mu\nu\rho\sigma}\Tr(G_{\mu\nu}G_{\rho\sigma})$。费米子部分：量纲 $3$ 的标量双线性型为 $\bar\psi\psi$、$\bar\psi\ii\gamma_5\psi$；量纲 $4$ 的为 $\bar\psi\Gamma^\mu D_\mu\psi$，$\Gamma^\mu\in\{\gamma^\mu,\gamma^\mu\gamma_5\}$（$\sigma^{\mu\nu}D_\mu D_\nu$ 的量纲为 $5$，四费米子算符的量纲为 $6$）；全导数 $\partial_\mu(\bar\psi\Gamma^\mu\psi)$ 把 $(\overline{D_\mu\psi})\Gamma^\mu\psi$ 化为已列出的形式，Hermite 性约束系数矩阵。第三步：分别对左、右手场作味空间的线性变换，使动能矩阵化为单位矩阵，再用双幺正变换把质量矩阵对角化为实而非负的对角阵。详细的讨论见\cite[第~15~章]{Weinberg1996}与\cite[第~15~章]{PeskinSchroeder1995}。
\end{proof}

\cref{thm:qcd-gauge-invariance}(b) 是“可重整的规范不变拉氏量本质上唯一”的经典层次版本：一旦规定规范群、夸克表示与可重整性（量纲 $\le4$），拉氏量就只剩下耦合常数 $g$、夸克质量与 $\theta$ 这几个参数。重整化之后这一结论仍然成立（\cref{thm:qcd-renormalizability}）：抵消项也只能是这些规范不变项，再加上规范固定部分的场与参数重整化。

\begin{remark}[$\theta$ 项]\label{rem:qcd-theta}
可以在\cref{eq:qcd-lagrangian}上加上
\[
\mathcal L_\theta=\frac{\theta g^2}{32\pi^2}\,G^a_{\mu\nu}\tilde G^{a\mu\nu}.
\]
由\cref{eq:qcd-chern-simons}，它是全导数；全导数项的顶点正比于流入顶点的总动量，由动量守恒而为零，所以 $\mathcal L_\theta$ 对微扰展开的任何一阶都没有贡献，只在非微扰层次上起作用。$G^a_{\mu\nu}\tilde G^{a\mu\nu}\propto\boldsymbol E^a\cdot\boldsymbol B^a$，故 $\mathcal L_\theta$ 破坏 P 与 T，而保持 C。强 CP 问题与瞬子不在本书范围之内；瞬子效应正比于 $\ee^{-8\pi^2/g^2}=\ee^{-2\pi/\as}$，对 $\as$ 的任何幂次都为零，它们将在第~1~章的全局假设\cref{asm:intro-af-coeff}的地位说明中再次出现。关于轴矢反常与 $\theta$ 参数，参见\cite[第~19~章]{PeskinSchroeder1995}。
\end{remark}

